\section{Uso de ferramentas de IA}
\textbf{IA como componente do experimento} (documentada na Metodologia): o bi-encoder \texttt{\EncName}, o cross-encoder \texttt{\CEName} e a LLM do RAG. \textbf{IA como apoio ao desenvolvimento:} usamos o assistente Claude (Anthropic) para escrever e depurar o código do projeto (índice invertido, BM25, correção ortográfica, métricas, scripts e testes), planejar e executar os experimentos, gerar tabelas e figuras, rascunhar o texto deste relatório e as análises por consulta a partir das evidências numéricas produzidas pelo código, e redigir as respostas das duas condições do experimento de RAG. Os integrantes revisaram o código, os resultados e o texto e são responsáveis por compreendê-los e defendê-los. A avaliação humana do RAG (notas de 0 a 2) cabe aos integrantes e não foi delegada à IA.
% CONFIRMAR antes de entregar: ajuste esta seção ao uso real de ferramentas de IA pelo grupo.
